% 张三慧V2第二十章《电磁感应》章末习题独立解析（扫描PDF205--208）。
% 公式统一LaTeX；几何方向和理想化边界需学术专家复核。
\section*{第20章\quad 电磁感应：独立解析}

\paragraph{习题 20.1}
长直导线在距轴 \(x\) 处
\(B(x)=\mu_0I/(2\pi x)\)；
由图20.22电流向上，
导线右边磁场\emph{入纸面}。
横杆向上运动，
\(\boldsymbol v\times\boldsymbol B\) 指向左端 \(a\)，
正电荷在 \(a\) 聚集。
杆沿 \(x\) 方向各处动生电势微元
\(d\mathcal E=vB(x)\,dx\)，所以
\[
 \boxed{|\mathcal E|=\int_d^{d+l}vB(x)\,dx
 =\frac{\mu_0Iv}{2\pi}\ln\frac{d+l}{d}}.
\]
代入 \(I=5\,\mathrm A,v=10\,\mathrm{m/s}\)、
\(d=.10\,\mathrm m,l=.20\,\mathrm m\)：
\[
 \boxed{|\mathcal E|\approx1.10\times10^{-5}\,\mathrm V},
 \qquad \boxed{\varphi_a>\varphi_b}.
\]

\paragraph{习题 20.2}
线圈半径 \(R=.12\,\mathrm m\)，
匝数 \(N=4000\)，
\(B=.5\,\mathrm G=5\times10^{-5}\,\mathrm T\)，
旋转频率 \(f=30\,\mathrm{s^{-1}}\)。
线圈磁链 \(N\Phi=NBA\cos(\omega t+\phi)\)，
\(\omega=2\pi f\)，
\(\mathcal E=-d(N\Phi)/dt\)。
因此
\[
 \boxed{\mathcal E_{\max}=NB\pi R^2(2\pi f)
 \approx1.71\,\mathrm V}.
\]
要达到最大值，线圈旋转轴
需与 \(\boldsymbol B\) 垂直，
使线圈法向扫过磁场方向；
在\emph{线圈平面平行于磁场}
（法向与磁场垂直、磁通瞬时为零）
时感应电动势的绝对值最大。
不可把 \(4\times10^3\) 匝误录成4匝。

\paragraph{习题 20.3}
穿过矩形线圈每匝的磁通
\[
 \Phi=\int_d^{d+a}\frac{\mu_0I}{2\pi x}L\,dx
 =\frac{\mu_0IL}{2\pi}\ln\frac{d+a}{d}.
\]
线圈以 \(v=dd/dt>0\) 远离直导线，
磁通减少，电势幅值
\[
 \boxed{|\mathcal E|
 =\frac{\mu_0NILv}{2\pi}
 \left(\frac1d-\frac1{d+a}\right)}.
\]
\(N=10^3,I=5\,\mathrm A,L=.20\,\mathrm m\)，
\(v=2\,\mathrm{m/s}\)，
\(d=a=.10\,\mathrm m\)，
\[
 \boxed{|\mathcal E|=2.00\times10^{-3}\,\mathrm V}.
\]
线圈上下边的动生电动势相消，
近侧与远侧竖直边的贡献不完全抵消。

\paragraph{习题 20.4}
图20.23中线圈固定，
电流
\(i(t)=5\sin(100\pi t)\,\mathrm A\)。
磁链为
\[
 \Psi(t)=\frac{\mu_0NL}{2\pi}
 \ln\frac{d+a}{d}\ i(t).
\]
故
\[
 \boxed{\mathcal E(t)
 =-\frac{\mu_0NL}{2\pi}
 \ln\frac{d+a}{d}\ \frac{di}{dt}}.
\]
由 \(di/dt=500\pi\cos(100\pi t)\)，
\(N=1000,L=.20,d=a=.10\)：
\[
 \boxed{\mathcal E(t)\approx
 -4.36\times10^{-2}\cos(100\pi t)\,\mathrm V}.
\]
负号依取的线圈环向正方向；
感生电势幅度约43.6mV。

\paragraph{习题 20.5}
圆盘区域 \(r\le R\) 有均匀随时间变化的
\(\boldsymbol B(t)\)，
图20.24直杆 \(ab\) 为低于圆心的水平弦。
圆心至弦垂距
\(h=\sqrt{R^2-L^2/4}\)，
三角区域 \(Oab\) 面积
\(S=Lh/2\)。
由法拉第定律对
\(O\to a\to b\to O\) 的闭路积分，
两径向边因为感生电场沿同心圆切向而
没有线积分，只有弦段贡献：
\[
 \boxed{|\mathcal E_{ab}|=
 \frac{L}{2}\sqrt{R^2-\frac{L^2}{4}}\,
 \left|\frac{dB}{dt}\right|}.
\]
图中磁场\emph{指向纸内且增强}，
按楞次定律感生电场逆时针，
在位于圆心下方的弦段从 \(a\) 向 \(b\)。
所以
\(\boxed{\varphi_b>\varphi_a}\)。
同半径两端即使径向距离相等，
也不能错误认为沿非保守感生电场弦段
的积分为零。

\paragraph{习题 20.6}
取题图竖直地磁场
\(B_z=.42\times10^{-4}\,\mathrm T\)，
两翼展长 \(l=12.705\,\mathrm m\)，
飞行速率
\(v\simeq1.70(330)=561\,\mathrm{m/s}\)
（配套辅导书使用声速330m/s）。
两翼尖横跨磁场动生电势
\[
 \boxed{|\mathcal E|=B_zlv
 =.42\times10^{-4}(12.705)(561)
 \approx0.30\,\mathrm V}.
\]
飞机从东向西飞、当地竖直磁场向下，
\(\boldsymbol v\times\boldsymbol B\)
由北指向南；
正电荷向南端翼尖移动，
故\(\boxed{\text{南翼尖电势高}}\)。
当地声速若取340m/s，
幅值约0.31V；需注明声速约定。

\paragraph{习题 20.7}
海水穿过垂直其速度方向的地磁竖直分量，
电极间距 \(l=200\,\mathrm m\)，
测得 \(|U|=7.0\times10^{-3}\,\mathrm V\)。
在理想均匀流条件下
\[
 \boxed{v=\frac{|U|}{Bl}
 =\frac{7.0\times10^{-3}}
 {(.70\times10^{-4})(200)}
 =0.50\,\mathrm{m/s}}.
\]
这一估算需电极间导电水流与磁场正交；
不同流场剖面会影响测得平均速度。

\paragraph{习题 20.8}
线圈面积 \(S=.10(.20)=.0200\,\mathrm{m^2}\)，
\(N=120,B=.020\,\mathrm T\)。
转动发电机感应电势幅值
\(\mathcal E_{\max}=NBS\,\omega\)。
令幅值12.0V，
\[
 \boxed{\omega=\frac{12.0}{120(.020)(.020)}
 =250\,\mathrm{rad/s}},
\quad
 \boxed{n_{\rm rpm}=\frac{60\omega}{2\pi}
 \approx2.39\times10^3\,\mathrm{r/min}}.
\]
配套书有时写旋转频率
\(f\approx40\,\mathrm{Hz}\)，与此一致。

\paragraph{习题 20.9}
线圈从磁通初值
\(\Psi_i=NBS\) 离开磁场至
\(\Psi_f=0\)，
法拉第定律
\(\mathcal E=-d\Psi/dt\)、
电流 \(i=\mathcal E/R\)。
总电量
\[
 q=\int i\,dt
 =-\frac1R\int_{\Psi_i}^{\Psi_f}d\Psi
 =\frac{NBS}{R}.
\]
因此
\[
 \boxed{B=\frac{qR}{NS}}.
\]
由于对\emph{磁链总变化}积分，
只要电阻保持常数且磁通确实从 \(NSB\) 变零，
最后的 \(q\) 不依赖移出运动是否匀速。

\paragraph{习题 20.10}
源图20.26中位于半径 \(r\) 的
小正方形区域，边长 \(a\ll r\)，
盘面厚度 \(b\)，电阻率 \(\rho\)。
当地盘速率 \(v=\omega r\)，
径向动生电动势
\(\mathcal E\simeq Bva=B\omega r a\)。
按配套辅导书图20.4使用近似：
小方块径向长度 \(a\) 的内回路电阻
\(R_{\rm patch}=\rho a/(ab)=\rho/b\)，
其余圆盘作为低阻回路。
故局部电流
\(I=\mathcal E/R_{\rm patch}
 =B\omega r ab/\rho\)。
垂直场中的切向制动力
\(F\simeq B I a\)，
阻力矩
\[
 \boxed{|\tau|=
 \frac{B^2a^2r^2b}{\rho}\,\omega}.
\]
阻力方向与 \(\omega\) 相反，
维持恒速所需机械功率
\[
 \boxed{P_{\rm mech}=|\tau|\omega
 =\frac{B^2a^2r^2b}{\rho}\omega^2}.
\]
磁制动功转成圆盘内欧姆热，
该量是图示局部低阻回流简化模型，
不是完整二维涡流边值问题的精确解。

\paragraph{习题 20.11}
轨道半径固定为 \(r\)，
带电量大小 \(e\)，
设轨道处的磁场 \(B_r(t)\)。
电子径向平衡
\(e vB_r=mv^2/r\)，即
\(mv=e r B_r\)；
保持 \(r\) 常数则
\[
 m\frac{dv}{dt}=er\frac{dB_r}{dt}.
\]
半径 \(r\) 圆周上的
非保守感生电场，由麦克斯韦—法拉第定律
\[
 2\pi r E_\theta
 =-\frac{d}{dt}
 \left(\pi r^2\overline B\right),
\qquad
 |E_\theta|=\frac r2
 \left|\frac{d\overline B}{dt}\right|.
\]
电子切向加速要求
\(m\,dv/dt=e|E_\theta|\)。
组合得
\[
 \boxed{\frac{dB_r}{dt}
 =\frac12\frac{d\overline B}{dt}}.
\]
若初始两场均零或
\(B_r=\overline B/2\) 的初值已成立，
积分得到
\(\boxed{B_r=\overline B/2}\)。
\emph{源条件限制：}
如果起始有预置磁通，
关系应含积分常数，
不能只由导数关系无条件推断绝对场值。

\paragraph{习题 20.12}
固定半径 \(r\) 的电子轨道
需向心力 \(F_c=mv^2/r\)。
速度增加 \(\Delta v\ll v\) 时，
\[
 \Delta F_c=
 \frac m r[(v+\Delta v)^2-v^2]
 \simeq\frac{2mv}{r}\Delta v.
\]
原有磁场满足 \(eBv=mv^2/r\)，
即 \(eB=mv/r\)。
当场也增加 \(\Delta B\)，
新的磁洛伦兹力的一级增量
\[
 \Delta F_B=e[(v+\Delta v)(B+\Delta B)-vB]
 \simeq eB\Delta v+ev\Delta B.
\]
要求保持半径，令
\(\Delta F_B=\Delta F_c\)，得
\[
 e v\Delta B=\frac{mv}{r}\Delta v
 \Longrightarrow
 \boxed{\Delta B=\frac{m}{er}\Delta v}.
\]
这正是习题20.11的加速器条件
\(B=mv/(er)\) 的微分形式。
速率接近光速时需改用相对论动量
而非 \(mv\)。

\paragraph{习题 20.13}
无限长近似的同轴均匀线圈组，
线圈1电流 \(I_1\) 在管内场
\(B_1=\mu_0N_1I_1/l\)。
每匝线圈2的磁通
\(\Phi_{21}=B_1S\)，
总磁链 \(N_2\Phi_{21}\)，
互感
\[
 \boxed{M_{21}=\frac{N_2\Phi_{21}}{I_1}
 =\frac{\mu_0N_1N_2S}{l}
 =\frac{\mu_0N_1N_2\pi R^2}{l}}.
\]
反向交换电流源同理得
\(\boxed{M_{12}=M_{21}}\)，
由磁场能交叉项的对称性也可解释。

\paragraph{习题 20.14}
大圆线圈 \(b\) 中心的磁感应强度
\(B_b=\mu_0N_bI_b/(2R_b)\)。
小线圈 \(a\) 面积远小于大圈，
磁场近似均匀，
\[
 \boxed{M=\frac{\mu_0N_aN_bS_a}{2R_b}}.
\]
代入
\(N_a=50,N_b=100,
S_a=4.0\times10^{-4}\,\mathrm{m^2},R_b=.20\,\mathrm m\)，
得
\[
 \boxed{M\approx6.28\times10^{-6}\,\mathrm H}.
\]
当 \(a\) 中电流以
\(dI_a/dt=-50\,\mathrm{A/s}\) 下降，
穿过圈 \(b\) 的\emph{每匝}磁通变化率
\[
 \boxed{\left|\frac{d\Phi_b}{dt}\right|
 =\frac{M}{N_b}\left|\frac{dI_a}{dt}\right|
 \approx3.14\times10^{-6}\,\mathrm{Wb/s}},
\]
线圈 \(b\) 的感应电动势幅值
\[
 \boxed{|\mathcal E_b|=M|dI_a/dt|
 \approx3.14\times10^{-4}\,\mathrm V}.
\]
\emph{磁链}变化率 \(N_b\,d\Phi_b/dt\)
比单匝磁通变化率大 \(N_b\) 倍，不能混报。

\paragraph{习题 20.15}
长直螺线管自感
\[
 L_s=\frac{\mu_0N^2A}{l},
\quad A=\pi r^2.
\]
取 \(r=.020\,\mathrm m,l=.300\,\mathrm m,N=1200\)，
\[
 \boxed{L_s=
 \frac{(4\pi\times10^{-7})(1200)^2\pi(.020)^2}{.300}
 \approx7.58\times10^{-3}\,\mathrm H}.
\]
电流变化率
\(3.0\times10^2\,\mathrm{A/s}\)，
\[
 \boxed{|\mathcal E_L|=L_s|dI/dt|
 \approx2.27\,\mathrm V}.
\]
方向总按楞次定律阻碍电流变化。
